% !TeX program = xelatex
% 核心检测识别程序代码框架 — 自然堆积赛道
% 编译: latexmk -xelatex framework.tex

\documentclass[UTF8,a4paper,landscape,10.5pt,fontset=fandol]{ctexart}
\usepackage[a4paper,landscape,top=1.2cm,bottom=1.2cm,left=1.35cm,right=1.35cm]{geometry}
\usepackage{tikz}
\usetikzlibrary{positioning,calc,arrows.meta}
\usepackage{booktabs}
\usepackage{array}
\usepackage{tabularx}
\usepackage{xcolor}
\usepackage{amsmath}
\usepackage{hyperref}
\pagestyle{empty}
\setlength{\parindent}{0pt}
\setlength{\tabcolsep}{4pt}
\renewcommand{\arraystretch}{1.16}
\hypersetup{hidelinks,pdftitle={核心检测识别程序代码框架 — 自然堆积赛道},pdfauthor={匿名}}

\definecolor{accentblue}{RGB}{40,91,143}
\definecolor{accentgreen}{RGB}{69,132,91}
\definecolor{accentorange}{RGB}{205,124,54}
\definecolor{accentpurple}{RGB}{113,83,142}
\definecolor{accentred}{RGB}{176,83,83}
\definecolor{accentgray}{RGB}{95,101,107}
\newcolumntype{Y}{>{\raggedright\arraybackslash}X}

\begin{document}

{\zihao{2}\bfseries 混凝土骨料颗粒智能筛分比拼}\hfill
{\zihao{4}自然堆积赛道}

\vspace{0.25em}
{\zihao{3}\bfseries 核心检测识别程序代码框架}
\hfill
{\small\textcolor{gray!80}{ImageGrains 2.0 / Cellpose-SAM + \texttt{aggregate\_screening}}}

\vspace{0.55em}

\begin{center}
\begin{tikzpicture}[
  font=\small,
  >=Latex,
  arr/.style={->,very thick,draw=accentgray},
  main/.style={rectangle,draw=accentgray!85,rounded corners=4pt,minimum height=1.45cm,text width=3.65cm,align=center,inner sep=5pt},
  small/.style={rectangle,draw=accentgray!85,rounded corners=4pt,minimum height=1.25cm,text width=3.45cm,align=center,inner sep=5pt},
  report/.style={rectangle,draw=accentgray!85,rounded corners=4pt,minimum height=1.35cm,text width=5.2cm,align=center,inner sep=5pt}
]

% 顶部主链
\node[main,fill=accentblue!7] (img) {\textbf{图像采集与尺度参考}\\\scriptsize 原始图像 + mm/px\\\texttt{app.py / scale.py}};
\node[main,fill=accentblue!7,right=0.70cm of img] (seg) {\textbf{实例分割}\\\scriptsize $I\rightarrow\{M_i\}$\\ImageGrains 2.0 / Cellpose-SAM};
\node[main,fill=accentgreen!9,right=0.70cm of seg] (geo) {\textbf{物理测量}\\\scriptsize $\{M_i\}\rightarrow\mathbf{x}_i$\\$a_i,b_i,A_i,P_i$};
\node[main,fill=accentorange!10,right=0.70cm of geo] (sieve) {\textbf{筛分映射}\\\scriptsize $d_i,w_i$\\\texttt{sieve\_equivalent.py}};

\draw[arr] (img) -- (seg);
\draw[arr] (seg) -- (geo);
\draw[arr] (geo) -- (sieve);

% 下游三支
\node[small,fill=accentpurple!9,below=1.05cm of geo,xshift=-3.20cm] (shape) {\textbf{投影形貌}\\\scriptsize $AR,C,S\rightarrow c_i$\\\texttt{morphology.py}};
\node[small,fill=accentred!8,below=1.05cm of geo,xshift=0.80cm] (anom) {\textbf{异常候选}\\\scriptsize 几何规则 $\rightarrow z_i$\\\texttt{anomaly.py}};
\node[small,fill=accentorange!10,below=1.05cm of sieve,xshift=2.00cm] (grade) {\textbf{质量代理级配}\\\scriptsize $D_{10},D_{50},D_{90},P_k$\\质量加权累计分布};

\draw[arr] (geo.south) -- (shape.north);
\draw[arr] (geo.south) -- (anom.north);
\draw[arr] (sieve) -- (grade);

% 汇总
\node[report,fill=accentgray!8,below=1.05cm of anom,xshift=0.80cm] (report) {\textbf{场景级报告与可视化}\\\scriptsize 逐颗粒表、级配曲线、形貌/异常汇总\\\texttt{report.py / visualization.py}};
\draw[arr] (shape.south) -- (report.north west);
\draw[arr] (anom.south) -- (report.north);
\draw[arr] (grade.south) -- (report.north east);

\end{tikzpicture}
\end{center}

\vspace{0.45em}

\begin{tabularx}{\linewidth}{>{\bfseries}p{4.2cm} Y Y}
\toprule
程序模块 & 主要职责 & 主要输出 \\
\midrule
\texttt{imagegrains} & 预训练实例分割与逐颗粒几何测量 & integer mask、长短轴、面积、周长等逐颗粒特征 \\
\texttt{scale.py} & 钢尺/既定分辨率的尺度估计与统一 & mm/px 及毫米级几何量 \\
\texttt{sieve\_equivalent.py} & 筛分等效粒径、质量代理与级配统计 & $d_i$、$w_i$、$D_{10}/D_{50}/D_{90}$、筛级占比 \\
\texttt{morphology.py} & 投影形貌分类 & $c_i$ \\
\texttt{anomaly.py} & 异常候选筛查 & $z_i$ \\
\texttt{report.py} / \texttt{visualization.py} & 汇总并导出场景级结果 & SceneSummary、CSV、表格、曲线与标注图 \\
\bottomrule
\end{tabularx}

\vspace{0.55em}
{\small
\textbf{当前固定口径：}筛分等效粒径采用 $\boldsymbol{\theta}=(1,0,0)$ 的短轴基线，质量代理采用 $w=d^3$。代码保留外部机械筛分数据可用时的校准接口，但该接口不参与当前报告结果计算。}

\end{document}
